% concatenated from expander_evolution_algebras.tex
%% Expander Evolution Algebras
%% Journal of Algebra submission
%% Compiled with: pdflatex expander_evolution_algebras.tex

\documentclass[12pt,a4paper]{amsart}

\usepackage{amsmath,amsthm,amssymb,amsfonts}
\usepackage{mathrsfs}
\usepackage{hyperref}
\usepackage{enumitem}
\usepackage{microtype}
\usepackage{doi}

% ------------------------------------------------------------------
% Theorem environments
% ------------------------------------------------------------------
\newtheorem{theorem}{Theorem}[section]
\newtheorem{proposition}[theorem]{Proposition}
\newtheorem{lemma}[theorem]{Lemma}
\newtheorem{corollary}[theorem]{Corollary}

\theoremstyle{definition}
\newtheorem{definition}[theorem]{Definition}
\newtheorem{example}[theorem]{Example}
\newtheorem{remark}[theorem]{Remark}
\newtheorem{notation}[theorem]{Notation}

\newtheorem{problem}{Open Problem}

% ------------------------------------------------------------------
% Macros
% ------------------------------------------------------------------
\DeclareMathOperator{\Span}{span}
\DeclareMathOperator{\Aut}{Aut}
\DeclareMathOperator{\Der}{Der}
\DeclareMathOperator{\val}{val}
\DeclareMathOperator{\diam}{diam}
\DeclareMathOperator{\supp}{supp}
\DeclareMathOperator{\Spec}{Spec}
\DeclareMathOperator{\tr}{tr}
\DeclareMathOperator{\rank}{rank}

\newcommand{\K}{\mathbb{K}}
\newcommand{\R}{\mathbb{R}}
\newcommand{\C}{\mathbb{C}}
\newcommand{\Z}{\mathbb{Z}}
\newcommand{\N}{\mathbb{N}}
\newcommand{\F}{\mathbb{F}}

% Graph notation
%\newcommand{\Gamma}{\varGamma}
\newcommand{\hh}{h}          % Cheeger / expansion constant
\newcommand{\EW}{E(W)}       % edge boundary of W

% Evolution algebra notation
\newcommand{\EA}{E}          % evolution algebra
\newcommand{\natbasis}{\mathcal{B}}  % natural basis
\newcommand{\strmat}{A}      % structural matrix
\newcommand{\evop}{\mathcal{L}} % evolution operator

% ------------------------------------------------------------------
\title{Expander Evolution Algebras}

\author{Piero Giacomelli}
\address{IT Department, Tenax Group \\ Via I Maggio, 226, 37020 Volargne (VR), Italy}
% \address{Department of Mathematics, [University Name] \\ [University Address]} 
\email{pgiacome@gmail.com}
\urladdr{https://github.com/PieroGiacomelli} % Added based on your research repositories

%\thanks{The author would like to thank...} % Optional: For acknowledgments or funding

\date{\today}

\keywords{Evolution algebra; expander graph; Cheeger constant;
          evolution operator; spectral gap; Ramanujan algebra;
          nonassociative algebra; Markov chain}
\subjclass[2020]{17D92, 05C50, 15A18, 60J10}

% ------------------------------------------------------------------
\begin{document}
% ------------------------------------------------------------------

\begin{abstract}
We introduce \emph{expander evolution algebras} (EEAs), a class of
nonassociative algebras defined over an arbitrary field $\K$ in which
the underlying undirected loopless graph of the algebra---in the sense
of Kowalski---is an expander graph in the classical sense of Cheeger.
Starting from the formal graph definition of Kowalski and the algebraic
framework of Tian, we establish a dictionary between combinatorial
expansion and algebraic structure: the Cheeger constant of the
associated graph governs connectivity, the subalgebra lattice, the
growth of the evolution sequence, and---over $\R$ and $\C$---the
spectral gap of the evolution operator.  Over a general field $\K$ we
prove that EEAs are always connected and simple (as evolution algebras),
carry no proper large evolution subalgebras, and that every generator of
a \emph{symmetric} EEA is algebraically persistent.  Over $\C$ we
obtain the sharp Alon--Boppana lower bound for the second eigenvalue of
the evolution operator, leading to the definition of
\emph{Ramanujan evolution algebras} as optimal expanders.  We also
construct families of EEAs from Cayley graphs of finite groups.
We close with open problems.
\end{abstract}

\maketitle

% ==================================================================
\section{Introduction}
% ==================================================================

Evolution algebras, introduced by Tian and Vojtechovsky \cite{TianVojtechovsky2006}
and developed systematically in the monograph \cite{Tian2008}, are
nonassociative, commutative, and non-power-associative algebras over a
field $\K$.  Their defining feature is a distinguished \emph{natural
basis} $\{e_1,\ldots,e_n\}$ in which distinct basis elements multiply
to zero, while $e_i^2 = \sum_k a_{ik} e_k$ for structural constants
$a_{ik}\in\K$.  Behind this rule lies varied structure: discrete dynamical systems, Markov chains,
non-Mendelian genetics, and graph combinatorics all admit natural
descriptions within this framework.

Expander graphs, in contrast, are finite graphs that are simultaneously
\emph{sparse} (bounded degree) and \emph{highly connected}
(large Cheeger constant).  Introduced implicitly by Pinsker \cite{Pinsker1973}
and formalized by Alon and Milman \cite{AlonMilman1985}, they have
become a central object in combinatorics, computer science, and number
theory; see Kowalski's comprehensive treatment \cite{Kowalski2019}.

The present paper asks: \emph{What algebraic properties does an
evolution algebra inherit when its associated graph is an expander?}
The answer turns out to connect expansion to connectivity, the subalgebra lattice,
evolution dynamics, and the spectral theory of the algebra.

\subsection*{Main contributions}

\begin{enumerate}[label=(\arabic*)]
\item We define the \emph{underlying undirected loopless graph}
      $\varGamma(\EA,\natbasis)$ of an evolution algebra $\EA$
      (Definition~\ref{def:assoc-graph}) following Kowalski's formal
      setup \cite[Definition~2.1.1]{Kowalski2019}, and use the
      Cheeger constant $\hh(\varGamma)$
      \cite[Definition~3.1.1]{Kowalski2019} to define
      \emph{Expander Evolution Algebras} (Definition~\ref{def:EEA}).

\item We prove that $\hh(\varGamma(\EA,\natbasis))>0$ if and only if
      $\EA$ is a connected evolution algebra
      (Theorem~\ref{thm:connectivity}), and that for nonsingular EEAs
      the isomorphism class of the associated graph is an algebraic
      invariant (Theorem~\ref{thm:basis-invariance}).

\item We prove a \emph{logarithmic diameter bound}
      (Theorem~\ref{thm:diameter}) and a \emph{support growth theorem}
      for the evolution sequence (Theorem~\ref{thm:support-growth}).

\item For \emph{symmetric} EEAs ($a_{ij}=a_{ji}$), we prove that every
      generator is algebraically persistent
      (Theorem~\ref{thm:persistency}), and that the algebra is simple
      (Theorem~\ref{thm:simplicity}).

\item In the $d$-regular case we prove a \emph{Cheeger inequality for
      the evolution operator} over $\R$
      (Theorem~\ref{thm:cheeger-evop}) and, over $\C$, the
      \emph{Alon--Boppana bound} (Theorem~\ref{thm:alon-boppana}),
      leading to the notion of a \emph{Ramanujan evolution algebra}
      (Definition~\ref{def:Ramanujan}).

\item We construct EEA families from \emph{Cayley graphs} of finite
      groups (Section~\ref{sec:cayley}) and from tensor products
      (Section~\ref{sec:tensor}).
\end{enumerate}

\subsection*{Organisation}
Section~\ref{sec:prelim} fixes notation and recalls the necessary
background on graphs, expanders, and evolution algebras.
Sections~\ref{sec:EEA}--\ref{sec:structure} contain the core algebraic
theory over a general field.  Section~\ref{sec:evop} studies the
evolution operator.  Section~\ref{sec:markov} specialises to Markov
(probabilistic) EEAs.  Section~\ref{sec:regular} develops the
$d$-regular theory and the spectral results.  Section~\ref{sec:complex}
gives the sharper results over $\C$.  Section~\ref{sec:constructions}
provides constructions.  Section~\ref{sec:open} states open problems.

% ==================================================================
\section{Preliminaries}\label{sec:prelim}
% ==================================================================

\subsection{Graphs in the sense of Kowalski}

We follow Kowalski's formal definition \cite[Definition~2.1.1]{Kowalski2019}
verbatim.

\begin{definition}[Graph]\label{def:graph}
A \emph{graph} $\varGamma$ is a triple $(V,E,\mathrm{ep})$ where $V$
and $E$ are sets (of \emph{vertices} and \emph{edges} respectively) and
\[
  \mathrm{ep}: E \longrightarrow V^{(2)}
\]
is the \emph{endpoint map}, where $V^{(2)}$ denotes the collection of
subsets of $V$ of cardinality $1$ or $2$.  If $\alpha\in E$ and
$\mathrm{ep}(\alpha)=\{x\}$ (a singleton), then $\alpha$ is a
\emph{loop} at $x$.  If $|\mathrm{ep}(\alpha)|=2$ the edge is
\emph{proper}.  A graph is \emph{simple} if it has no loops and no two
distinct edges share the same pair of endpoints.  The \emph{valency} (or
\emph{degree}) of a vertex $x$ is $\val(x)=|\{\alpha\in E : x\in
\mathrm{ep}(\alpha)\}|$.  A graph is \emph{$d$-regular} if $\val(x)=d$
for every $x\in V$.  We write $|\varGamma|=|V|$ for the \emph{size} of
$\varGamma$.
\end{definition}

Throughout this paper, all graphs are assumed to be \emph{finite,
simple} (no loops, no multiple edges), and \emph{connected} unless
otherwise stated.

\begin{definition}[Adjacency matrix]\label{def:adj}
Let $\varGamma=(V,E,\mathrm{ep})$ be a finite graph with
$V=\{v_1,\ldots,v_n\}$.  The \emph{adjacency matrix} of $\varGamma$ is
the symmetric $\{0,1\}$-matrix $A_\varGamma=(a_{ij})$ where
$a_{ij}=|\{\alpha\in E : \mathrm{ep}(\alpha)=\{v_i,v_j\}\}|$.  For a
simple graph, $a_{ij}\in\{0,1\}$ and $a_{ii}=0$.
\end{definition}

\subsection{Expansion constant and expander graphs}

\begin{definition}[Cheeger constant]\label{def:cheeger}
Let $\varGamma=(V,E,\mathrm{ep})$ be a finite graph.  For a subset
$W\subseteq V$ let
\[
  E(W) \;=\; \{\alpha\in E : |\mathrm{ep}(\alpha)\cap W|=1\}
\]
be the \emph{edge boundary} of $W$.  The \emph{Cheeger constant} (or
\emph{expansion constant}) of $\varGamma$ is
\[
  \hh(\varGamma)
  \;=\;
  \min\!\left\{
    \frac{|E(W)|}{|W|} \;\Big|\;
    \emptyset \neq W \subsetneq V,\; |W|\le\tfrac{1}{2}|\varGamma|
  \right\},
\]
with the convention $\hh(\varGamma)=+\infty$ if $|V|\le 1$.
\end{definition}

The following basic facts are due to Kowalski
\cite[Proposition~3.1.2 and Lemma~3.1.4]{Kowalski2019}.

\begin{proposition}[Basic properties of $\hh$]\label{prop:cheeger-basic}
Let $\varGamma$ be a finite graph with $|V|\ge 2$.
\begin{enumerate}[label=\normalfont(\roman*)]
  \item $\hh(\varGamma)>0$ if and only if $\varGamma$ is connected.
  \item For any $W\subseteq V$ with $|W|=\delta|V|$, $0<\delta\le\frac12$,
        one must remove at least $\delta\,\hh(\varGamma)\,|V|$ edges to
        disconnect $W$ from its complement.
  \item If $\varGamma$ is connected then
        $\dfrac{2}{|V|}\le\hh(\varGamma)\le\min_{x\in V}\val(x)$.
\end{enumerate}
\end{proposition}

\begin{definition}[Expander graph]\label{def:expander}
A finite connected graph $\varGamma$ is called an
\emph{$h$-expander} for a constant $h>0$ if
$\hh(\varGamma)\ge h$.
A family $(\varGamma_m)_{m\ge 1}$ of finite connected graphs is an
\emph{expander family} if:
\begin{enumerate}[label=\normalfont(\roman*)]
  \item $|\varGamma_m|\to\infty$;
  \item the degrees are uniformly bounded: $\sup_{m}\max_{x\in V_m}\val(x)<\infty$;
  \item the expansion constants are uniformly bounded below:
        $\inf_m \hh(\varGamma_m)>0$.
\end{enumerate}
\end{definition}

\begin{theorem}[Diameter bound {\cite[Corollary~3.1.10]{Kowalski2019}}]
\label{thm:diam-expander}
Let $\varGamma$ be a finite connected graph with
$\hh(\varGamma)\ge h>0$ and $\max_{x}\val(x)\le d$.  Then
\[
  \diam(\varGamma) \;\le\;
  \frac{2d}{h}\,\log|\varGamma|  \;+\;1.
\]
\end{theorem}

\subsection{Evolution algebras}\label{sub:ea}

We follow Tian's monograph \cite{Tian2008} for the algebraic
definitions.

\begin{definition}[Evolution algebra]\label{def:ea}
Let $\K$ be a field.  An \emph{evolution algebra} over $\K$ is an
algebra $(\EA,\cdot)$ together with a countable \emph{natural basis}
$\natbasis=\{e_i\}_{i\in\Lambda}$ such that
\[
  e_i \cdot e_j = 0 \quad (i\ne j),
  \qquad
  e_i \cdot e_i = \sum_{k\in\Lambda} a_{ik}\,e_k \quad
  \text{for all }i\in\Lambda,
\]
where $a_{ik}\in\K$ and only finitely many $a_{ik}$ are nonzero for
each fixed $i$.  The scalars $(a_{ik})$ are the
\emph{structural constants} of $\EA$ with respect to $\natbasis$, and
the matrix $\strmat=(a_{ik})$ is the \emph{structural matrix}.  We
call $\dim\EA=|\Lambda|$ the \emph{dimension} of $\EA$.
\end{definition}

By Tian \cite[Theorem~1]{Tian2008}, when $\Lambda$ is finite the
natural basis $\natbasis$ is a linear basis for $\EA$.  We restrict to
the \emph{finite-dimensional} case throughout:
$n=\dim\EA<\infty$, $\Lambda=\{1,\ldots,n\}$.

\begin{remark}
Every evolution algebra is \emph{commutative} and \emph{flexible}, but
\emph{not} associative or power-associative in general
\cite[Corollary~1]{Tian2008}.
\end{remark}

\begin{definition}[Evolution subalgebra and ideal]\label{def:subalg}
Let $\EA$ be an evolution algebra with natural basis $\natbasis=
\{e_1,\ldots,e_n\}$.
\begin{enumerate}[label=\normalfont(\roman*)]
  \item A subspace $\EA'\subseteq\EA$ is an \emph{evolution subalgebra}
        if it has a natural basis $\{e_i\}_{i\in S}$ that extends to a
        natural basis of $\EA$.
  \item Every evolution subalgebra is automatically a two-sided
        \emph{evolution ideal}
        \cite[Proposition~2]{Tian2008}.
  \item $\EA$ is \emph{connected} if it cannot be written as
        $\EA=\EA_1\oplus\EA_2$ (direct sum of proper evolution subalgebras).
  \item $\EA$ is \emph{simple} if it has no proper evolution ideal.
\end{enumerate}
\end{definition}

\begin{definition}[Principal and plenary powers]\label{def:powers}
Let $x\in\EA$.  The \emph{principal powers} are $x^1=x$, $x^m=x^{m-1}\cdot x$ for
$m\ge 2$.  The \emph{plenary powers} are $x^{[0]}=x$, $x^{[k]}=
x^{[k-1]}\cdot x^{[k-1]}$ for $k\ge 1$.
\end{definition}

\begin{definition}[Algebraic persistency and transiency]\label{def:persist}
A generator $e_i$ of $\EA$ is \emph{algebraically persistent} if $e_i$
occurs (has nonzero coefficient) in $e_i^{[k]}$ for all $k\ge 0$.
Otherwise, $e_i$ is \emph{algebraically transient}.  We say $\EA$ is
\emph{algebraically persistent} if every generator is algebraically
persistent.
\end{definition}

\begin{definition}[Occurrence relation]\label{def:occur}
For generators $e_i,e_j$ of $\EA$, write $e_i\prec e_j$ if $e_i$
occurs in $e_j^2=\sum_k a_{jk} e_k$, i.e., $a_{ji}\ne 0$.  We extend
this to plenary powers: $e_i\prec^{[k]} e_j$ if $e_i$ occurs in
$e_j^{[k]}$.
\end{definition}

\subsection{The directed graph of an evolution algebra}

Elduque and Labra \cite{ElduqueLabra2013} associate a weighted directed
graph to any evolution algebra.

\begin{definition}[Weighted digraph]\label{def:digraph}
Let $\EA$ be an evolution algebra with natural basis $\natbasis=
\{e_1,\ldots,e_n\}$ and structural matrix $\strmat=(a_{ij})$.  The
\emph{weighted digraph} $D(\EA,\natbasis)$ has vertex set $\{1,\ldots,n\}$
and a directed edge $i\to j$ with weight $a_{ij}$ whenever
$a_{ij}\ne 0$ (and $i\ne j$).
\end{definition}

Key facts from Elduque--Labra \cite{ElduqueLabra2013}:
\begin{itemize}
  \item $\EA$ is nilpotent if and only if $D(\EA,\natbasis)$ has no
        directed cycle.
  \item $\Aut(\EA)$ is finite whenever $\EA^2=\EA$.
\end{itemize}

% ==================================================================
\section{Expander Evolution Algebras: Definition and First Properties}
\label{sec:EEA}
% ==================================================================

\subsection{The underlying undirected graph}

Since Kowalski's expander theory is developed for \emph{undirected} graphs, and the
structural matrix $\strmat$ need not be symmetric over a general field, we pass to the
symmetrization.

\begin{definition}[Underlying undirected loopless graph]\label{def:assoc-graph}
Let $\EA$ be a finite-dimensional evolution algebra over $\K$ with
natural basis $\natbasis=\{e_1,\ldots,e_n\}$ and structural matrix
$\strmat=(a_{ij})$.  The \emph{underlying undirected loopless graph}
$\varGamma(\EA,\natbasis)$ is the simple graph with:
\begin{itemize}
  \item vertex set $V=\{1,\ldots,n\}$;
  \item edge set $E_\varGamma=\bigl\{\{i,j\}: i\ne j,\;
                                    a_{ij}\ne 0 \text{ or } a_{ji}\ne 0\bigr\}$;
  \item endpoint map $\mathrm{ep}(\{i,j\})=\{i,j\}$.
\end{itemize}
In the notation of Definition~\ref{def:graph}, $\varGamma(\EA,\natbasis)
=(V,E_\varGamma,\mathrm{ep})$.
\end{definition}

\begin{remark}\label{rem:symm}
$\varGamma(\EA,\natbasis)$ is the \emph{symmetrisation} of the directed
graph $D(\EA,\natbasis)$: the edge $\{i,j\}$ exists in
$\varGamma(\EA,\natbasis)$ if and only if at least one of $i\to j$ or
$j\to i$ is present in $D(\EA,\natbasis)$.  The diagonal entries
$a_{ii}$ do not contribute edges to $\varGamma(\EA,\natbasis)$, since we
require loops to be excluded.
\end{remark}

\begin{remark}\label{rem:field}
The definition of $\varGamma(\EA,\natbasis)$ only uses the
\emph{support pattern} of $\strmat$---which entries are zero and which
are not---not their specific values.  Thus $\varGamma(\EA,\natbasis)$
is independent of the particular nonzero values of $a_{ij}$.
\end{remark}

\begin{definition}[Expansion constant of an evolution algebra]\label{def:ea-cheeger}
Let $\EA$ be a finite-dimensional evolution algebra with natural basis
$\natbasis$.  The \emph{expansion constant of $\EA$ with respect to
$\natbasis$} is
\[
  \hh(\EA,\natbasis) \;:=\; \hh\!\bigl(\varGamma(\EA,\natbasis)\bigr),
\]
where the right-hand side is the Cheeger constant of Definition
\ref{def:cheeger}.
\end{definition}

\begin{definition}[Expander Evolution Algebra]\label{def:EEA}
A finite-dimensional evolution algebra $\EA$ over $\K$ with natural
basis $\natbasis$ is called an \emph{$h$-Expander Evolution Algebra}
($h$-EEA) if
\[
  \hh(\EA,\natbasis) \;\ge\; h \;>\; 0.
\]
We call $h$ the \emph{expansion constant} of $\EA$.  A family
$(\EA_m,\natbasis_m)_{m\ge 1}$ with $\dim\EA_m\to\infty$ is an
\emph{EEA family} if $\inf_m \hh(\EA_m,\natbasis_m)>0$.
\end{definition}

\subsection{Connectivity and simplicity}\label{sub:connectivity}

\begin{theorem}[Connectivity]\label{thm:connectivity}
Let $\EA$ be a finite-dimensional evolution algebra over $\K$ with
natural basis $\natbasis=\{e_1,\ldots,e_n\}$.  The following are
equivalent:
\begin{enumerate}[label=\normalfont(\roman*)]
  \item $\EA$ is connected (Definition~\ref{def:subalg}(iii)).
  \item $\varGamma(\EA,\natbasis)$ is a connected graph.
  \item $\hh(\EA,\natbasis)>0$.
\end{enumerate}
\end{theorem}

\begin{proof}
The equivalence (ii)$\Leftrightarrow$(iii) is
Proposition~\ref{prop:cheeger-basic}(i).

\medskip
\noindent
\textbf{(i)$\Rightarrow$(ii).}
Suppose $\varGamma(\EA,\natbasis)$ is disconnected; let $S\sqcup
S^c=\{1,\ldots,n\}$ be a partition into two non-empty sets with no
edges between them.  This means: $a_{ij}=0$ and $a_{ji}=0$ whenever
$i\in S$, $j\in S^c$.  Let $\EA_1=\Span_\K\{e_i:i\in S\}$ and
$\EA_2=\Span_\K\{e_j:j\in S^c\}$.  For $i\in S$:
\[
  e_i^2 \;=\; \sum_k a_{ik}\,e_k \;=\; \sum_{k\in S} a_{ik}\,e_k \;\in\; \EA_1,
\]
because $a_{ij}=0$ for all $j\in S^c$.  Hence $\EA_1$ is an evolution
subalgebra.  By the same argument $\EA_2$ is an evolution subalgebra.
Since the basis of $\EA_1$ extends (with $\natbasis|_{S^c}$) to a
natural basis of $\EA$, the sum is direct: $\EA=\EA_1\oplus\EA_2$,
contradicting connectivity of $\EA$.

\medskip
\noindent
\textbf{(ii)$\Rightarrow$(i).}
Suppose $\EA=\EA_1\oplus\EA_2$ with $\EA_k$ proper evolution
subalgebras, generated respectively by $\{e_i:i\in S\}$ and
$\{e_j:j\in S^c\}$ with $S\sqcup S^c=\{1,\ldots,n\}$.  Since $\EA_1$
is an evolution subalgebra, for every $i\in S$:
$e_i^2\in\EA_1$, so $a_{ij}=0$ for all $j\in S^c$.  Symmetrically
$a_{ji}=0$ for all $j\in S^c$, $i\in S$.  Hence there are no edges
between $S$ and $S^c$ in $\varGamma(\EA,\natbasis)$, contradicting
connectivity.
\end{proof}

\begin{corollary}\label{cor:EEA-connected}
Every EEA is connected.
\end{corollary}

\begin{theorem}[Simplicity of symmetric EEAs]\label{thm:simplicity}
Let $\EA$ be a finite-dimensional evolution algebra over $\K$ with
natural basis $\natbasis=\{e_1,\ldots,e_n\}$.  Suppose $\EA$ is
\emph{symmetric}, meaning $a_{ij}=a_{ji}$ for all $i\ne j$.  If $\EA$
is an $h$-EEA for some $h>0$, then $\EA$ is simple.
\end{theorem}

\begin{proof}
By Definition~\ref{def:subalg}(iii) and
Tian~\cite[Corollary~2]{Tian2008}, simplicity is equivalent to
connectivity for evolution algebras over $\K$ (since every evolution
subalgebra is an ideal and the algebra is connected if and only if it
has no proper ideal).  The result then follows immediately from
Theorem~\ref{thm:connectivity}.
\end{proof}

\begin{remark}
For non-symmetric EEAs, simplicity and connectivity may differ.  We
will return to this point in Section~\ref{sub:markov-simple}.
\end{remark}

\subsection{Basis invariance}\label{sub:basis-inv}

The graph $\varGamma(\EA,\natbasis)$ depends a priori on the choice of
natural basis.  We now show that for \emph{nonsingular} evolution
algebras, the isomorphism class of $\varGamma(\EA,\natbasis)$ is a
genuine algebraic invariant.

\begin{theorem}[Basis invariance]\label{thm:basis-invariance}
Let $\EA$ be a finite-dimensional evolution algebra over $\K$ with
structural matrix $\strmat$ of rank $n$ (nonsingular).  Then any two
natural bases $\natbasis$ and $\natbasis'$ of $\EA$ give rise to
isomorphic underlying graphs:
\[
  \varGamma(\EA,\natbasis) \;\cong\; \varGamma(\EA,\natbasis').
\]
\end{theorem}

\begin{proof}
By Casas, Ladra, Omirov, and Rozikov \cite{CasasLadraOmirovRozikov2010},
the automorphisms of a nonsingular evolution algebra (i.e., the maps
that send a natural basis to another natural basis) are precisely the
compositions of:
\begin{enumerate}[label=\normalfont(\roman*)]
  \item a permutation $\sigma$ of the basis elements: $e_i\mapsto e_{\sigma(i)}$;
  \item scalar rescalings $e_i\mapsto\lambda_i e_i$ with
        $\lambda_i\ne 0$.
\end{enumerate}
Under a permutation $\sigma$, the structural matrix transforms as
$a'_{ij}=a_{\sigma^{-1}(i),\sigma^{-1}(j)}$, which gives the graph
$\varGamma(\EA,\natbasis')$ isomorphic to $\varGamma(\EA,\natbasis)$
via the vertex permutation $\sigma$.

Under a scalar rescaling $e_i\mapsto\lambda_i e_i$, the new structural
constants are $a'_{ij}=a_{ij}\lambda_j/\lambda_i^2$.  Since $\K$ is a
field and all $\lambda_i\ne 0$, we have $a'_{ij}\ne 0$ if and only if
$a_{ij}\ne 0$.  Hence the support pattern of $\strmat'$ equals that of
$\strmat$, so $\varGamma(\EA,\natbasis')=\varGamma(\EA,\natbasis)$
(same graph, not just isomorphic).
\end{proof}

\begin{corollary}\label{cor:EEA-invariant}
The property of being an $h$-EEA is an algebraic invariant of
nonsingular evolution algebras.
\end{corollary}

\begin{proof}
We argue in two steps.  First we establish that, for a fixed nonsingular
evolution algebra, the value $\hh(\EA,\natbasis)$ does not depend on the
chosen natural basis; we then transport this result along an arbitrary
isomorphism of nonsingular evolution algebras.

\medskip
\noindent\textbf{Step~1 (basis independence).}
Let $\natbasis$ and $\natbasis'$ be two natural bases of a nonsingular
evolution algebra $\EA$.  By
Theorem~\ref{thm:basis-invariance}, the graphs
$\varGamma(\EA,\natbasis)$ and $\varGamma(\EA,\natbasis')$ are
isomorphic as simple graphs.  The Cheeger constant of
Definition~\ref{def:cheeger} is defined entirely in terms of the
cardinalities $|W|$ and $|E(W)|$ for vertex subsets $W$, quantities
that are preserved by any bijection of vertices respecting the edge
set.  Hence $\hh(\cdot)$ is a graph isomorphism invariant, and
\[
  \hh(\EA,\natbasis)
  \;=\;\hh\bigl(\varGamma(\EA,\natbasis)\bigr)
  \;=\;\hh\bigl(\varGamma(\EA,\natbasis')\bigr)
  \;=\;\hh(\EA,\natbasis').
\]
In particular the condition $\hh(\EA,\natbasis)\ge h$ holds for one
natural basis if and only if it holds for every natural basis.

\medskip
\noindent\textbf{Step~2 (isomorphism invariance).}
Let $\varphi\colon\EA\to\EA'$ be an isomorphism of nonsingular
evolution algebras, and let $\natbasis=\{e_1,\ldots,e_n\}$ be a natural
basis of $\EA$ with structural constants $a_{ij}$.  Because $\varphi$
is bijective, linear, and multiplicative,
\[
  \varphi(e_i)\,\varphi(e_j)
  \;=\;\varphi(e_ie_j)
  \;=\;0
  \qquad (i\ne j),
\]
so $\varphi(\natbasis)=\{\varphi(e_1),\ldots,\varphi(e_n)\}$ is a
natural basis of $\EA'$.  Moreover
\[
  \varphi(e_i)^{\,2}
  \;=\;\varphi(e_i^{\,2})
  \;=\;\varphi\!\Bigl(\sum_{k=1}^{n}a_{ik}e_k\Bigr)
  \;=\;\sum_{k=1}^{n}a_{ik}\,\varphi(e_k),
\]
so the structural constants of $\EA'$ in the basis $\varphi(\natbasis)$
coincide with those of $\EA$ in $\natbasis$.  By
Definition~\ref{def:assoc-graph} this yields the equality of graphs
\[
  \varGamma\bigl(\EA',\varphi(\natbasis)\bigr)
  \;=\;
  \varGamma(\EA,\natbasis),
\]
and therefore $\hh(\EA',\varphi(\natbasis))=\hh(\EA,\natbasis)$.

\medskip
\noindent\textbf{Conclusion.}
Combining Step~1 (applied to $\EA'$) with Step~2, we obtain
\[
  \hh(\EA',\natbasis'')
  \;=\;\hh\bigl(\EA',\varphi(\natbasis)\bigr)
  \;=\;\hh(\EA,\natbasis)
\]
for every natural basis $\natbasis''$ of $\EA'$.  Consequently $\EA$
satisfies $\hh(\EA,\natbasis)\ge h$ if and only if $\EA'$ satisfies
$\hh(\EA',\natbasis'')\ge h$, which is precisely the assertion that
the $h$-EEA property is invariant under isomorphism of nonsingular
evolution algebras.
\end{proof}

% ==================================================================
\section{Structural Properties of EEAs}\label{sec:structure}
% ==================================================================

Section~\ref{sec:structure} develops the structural consequences of
the expansion hypothesis introduced in
Section~\ref{sec:EEA}.  The guiding question is the following: what
does it mean, \emph{algebraically}, for an evolution algebra to satisfy
a uniform isoperimetric inequality on its underlying undirected graph?
The expansion condition $\hh(\EA,\natbasis)\ge h>0$ is intrinsically
combinatorial; it bounds, from below, the ratio $|E(W)|/|W|$ over every
admissible subset $W$ of natural-basis indices.  In an evolution
algebra, however, subsets of indices correspond to evolution
subalgebras, plenary powers act on the basis through neighbourhood
enlargement, and the persistency dichotomy of Tian
\cite{Tian2008} classifies generators according to their long-time
return behaviour.  The aim of this section is to transport the
combinatorial content of the expansion inequality into each of these
algebraic settings, producing four structural conclusions: a
quantitative impossibility of small almost-isolated subalgebras and of
small direct summands, a logarithmic diameter bound, a near-exponential
support growth law for plenary powers, and a triviality theorem for
the hierarchical decomposition in the symmetric case.

\medskip
The first observation is that the very definition of $\hh(\EA,
\natbasis)$ already constrains the lattice of evolution subalgebras
generated by subsets of the natural basis.  The minimum defining the
expansion constant ranges over every nonempty
$S\subseteq\{1,\ldots,n\}$ with $|S|\le n/2$; consequently, for every
evolution subalgebra $\EA'=\Span_{\K}\{e_i:i\in S\}$ of dimension
$m\le n/2$, the boundary $E(S)$ in $\varGamma(\EA,\natbasis)$ contains
at least $h\cdot m$ edges.  Each such edge is a pair $\{i,j\}$ with
$i\in S$ and $j\notin S$, and at least one of the structural constants
$a_{ij}$ or $a_{ji}$ is nonzero.  Two consequences are immediate.
First, no proper evolution subalgebra of dimension at most $n/2$ can
be \emph{isolated} from its complement: it must be linked to it
through at least $hm$ nonzero structural constants.  Second, the
algebra cannot decompose as an internal direct sum $\EA=\EA_1\oplus
\EA_2$ with $\dim\EA_1\le n/2$, since such a decomposition would force
the boundary edge set to be empty.  These two facts should be read as
structural rigidity statements: the expansion hypothesis prevents the
appearance of an algebraically detached fragment of size up to half
the dimension, and they follow at once from
Definition~\ref{def:EEA}.  We will not record them as separate
theorems; they enter the rest of the section implicitly, through the
expansion inequality applied to the supports of plenary powers.

\medskip
The first formal result of the section is a logarithmic diameter
bound (Theorem~\ref{thm:diameter}).  Applying Kowalski's diameter
inequality for expander graphs to $\varGamma(\EA,\natbasis)$ yields
$\diam(\varGamma)\le (2d/h)\log n+1$, where $d$ is an upper bound for
the maximal valence.  The algebraic content of this bound becomes
transparent through the notion of \emph{algebraic distance}: the
distance between two generators $e_i$ and $e_j$ in $\varGamma(\EA,
\natbasis)$ coincides with the smallest $k$ such that $e_j$ occurs in
the $k$-th plenary power $e_i^{[k]}$, and equivalently with the
smallest $k$ for which $e_j$ appears in the support of $\evop^k(e_i)$.
The diameter bound therefore states that, in any EEA family with
bounded valence, every pair of generators is connected through a chain
of nonzero structural couplings of length $O(\log n)$: the number of
evolution steps separating any two generators grows only
logarithmically in the dimension of the algebra.

\medskip
The diameter bound is then sharpened by the support growth theorem
(Theorem~\ref{thm:support-growth}), which controls not just the
existence of long-distance couplings but also the rate at which the
support of plenary powers expands.  For symmetric EEAs the theorem
shows that $|S_k|\ge\min(n,(1+h/d)^k)$, where
$S_k=\supp(e_i^{[k]})$.  The proof exploits the elementary observation
that, while $|S_k|\le n/2$, the expansion inequality applied to $S_k$
produces at least $h|S_k|$ outgoing edges, each contributing a new
index to $S_{k+1}$ at a rate controlled by the maximal valence.  The
support therefore grows at least geometrically until it covers more
than half of the basis, and a single further step then completes the
cover.  Corollary~\ref{cor:cover-time} records the resulting
logarithmic cover time: every plenary power $e_i^{[k]}$ has full
support once $k$ exceeds $O((d/h)\log n)$.  The same mechanism
explains how a single generator suffices to reach the entire algebra
under iterated squaring, in a number of steps that scales with $\log
n$ rather than $n$, in striking contrast to the worst-case behaviour
admitted by general finite-dimensional evolution algebras.

\medskip
The third group of results concerns the symmetric case
$a_{ij}=a_{ji}$.  Theorem~\ref{thm:persistency} shows that every
generator of a symmetric EEA is \emph{algebraically persistent} in the
sense of Tian: the generator $e_i$ occurs in $e_i^{[k]}$ for all
$k\ge 0$.  The argument has two ingredients.  First, symmetry forces
every undirected edge of $\varGamma(\EA,\natbasis)$ to lift to a pair
of opposite directed edges in $D(\EA,\natbasis)$, so that every
neighbour of $i$ contributes a copy of $e_i$ when squared.  Second,
the support growth theorem ensures that, after at most $O((d/h)\log
n)$ steps, the support of $e_i^{[k]}$ is the full index set, and an
explicit local computation handles the transient regime $k<k_0$.  The
conclusion is that no generator of a symmetric EEA is transient: the
long-time behaviour of plenary powers is uniformly recurrent across
the basis.  This persistency phenomenon is sharp: the remark following
Theorem~\ref{thm:persistency} shows that symmetry is necessary, and
the non-symmetric case admits expander graphs whose evolution algebra
contains transient generators.

\medskip
The section closes with the hierarchical structure result
(Theorem~\ref{thm:hierarchy}): a symmetric EEA has trivial Tian
hierarchy, in the sense that its only level is the bottom level
containing all generators.  This is the structural shadow of universal
persistency.  In Tian's framework the levels of the hierarchy are
obtained by successively peeling off transient generators, and in the
symmetric expander setting there are none to peel off.  Combined with
Theorem~\ref{thm:connectivity} and Corollary~\ref{cor:EEA-connected},
which already secured connectivity, and with
Theorem~\ref{thm:simplicity}, which gave simplicity, this reduces the
structural picture of a symmetric EEA to the simplest possible form: a
single connected, simple, persistent block of dimension $n$ in which
every generator participates in the asymptotic dynamics.  Section
\ref{sec:evop} will then refine this picture by passing to the
spectrum of the evolution operator, where the expansion constant
reappears as a quantitative gap in the spectrum of $\strmat$.

\subsection{Diameter bound and evolution growth}

\begin{theorem}[Diameter bound]\label{thm:diameter}
Let $\EA$ be an $h$-EEA over $\K$ with natural basis $\natbasis$,
$n=\dim\EA$, and $\max_i\val(i)\le d$ in $\varGamma(\EA,\natbasis)$.
Then
\[
  \diam\!\bigl(\varGamma(\EA,\natbasis)\bigr)
  \;\le\;
  \frac{2d}{h}\,\log n \;+\; 1.
\]
In particular, for fixed $h>0$ and $d<\infty$, the diameter is
$O(\log n)$.
\end{theorem}

\begin{proof}
Direct application of Theorem~\ref{thm:diam-expander} to
$\varGamma(\EA,\natbasis)$.
\end{proof}

To make this algebraically meaningful, recall that the \emph{algebraic
distance} between generators $e_i$ and $e_j$ is the length of the
shortest path in $\varGamma(\EA,\natbasis)$ connecting $i$ and $j$,
which equals the smallest $k$ such that $e_j$ occurs in
$(e_i^2)^{[k]}$ (or in $e_i^{[k+1]}$ in the symmetric case).

\begin{definition}[Support of an element]\label{def:support}
For $x=\sum_k \alpha_k e_k\in\EA$, the \emph{support} of $x$ is
$\supp(x)=\{k : \alpha_k\ne 0\}\subseteq\{1,\ldots,n\}$.
\end{definition}

\begin{theorem}[Support growth]\label{thm:support-growth}
Let $\EA$ be a symmetric $h$-EEA over $\K$ with $\max_i\val(i)\le d$.
For any generator $e_i$, let $S_k=\supp(e_i^{[k]})$ be the support of
its $k$-th plenary power.  Then:
\begin{enumerate}[label=\normalfont(\roman*)]
  \item $|S_k|\ge\min\!\left(n,\; \left(1+\tfrac{h}{d}\right)^k\right)$
        for all $k\ge 0$.
  \item There exists $k_0\le\frac{d}{h}\log n$ such that
        $|S_{k_0}|=n$.
\end{enumerate}
\end{theorem}

\begin{proof}
Since $\EA$ is symmetric ($a_{ij}=a_{ji}$), the plenary power
$e_i^{[k+1]}=(e_i^{[k]})^2$.  We have
\[
  e_i^{[k+1]} \;=\; \left(\sum_{j\in S_k}\alpha_j^{(k)} e_j\right)^2
  \;=\; \sum_{j\in S_k} (\alpha_j^{(k)})^2\, e_j^2
  \;=\; \sum_{j\in S_k} (\alpha_j^{(k)})^2 \sum_\ell a_{j\ell}\,e_\ell.
\]
Hence $S_{k+1}\supseteq\bigcup_{j\in S_k} \{l : a_{j\ell}\ne 0\}
\supseteq S_k\cup N(S_k)$, where $N(S_k)$ is the neighbourhood of
$S_k$ in $\varGamma(\EA,\natbasis)$.

For part (i): as long as $|S_k|\le n/2$, the expansion condition gives
$|E(S_k)|\ge h|S_k|$.  Since each vertex in $S_k$ has degree at most
$d$, the number of new neighbours $|N(S_k)\setminus S_k|$ satisfies
\[
  |N(S_k)\setminus S_k| \;\ge\; \frac{|E(S_k)|}{d}\cdot
  \frac{|N(S_k)\setminus S_k|}{|N(S_k)\setminus S_k|}
  \;\ge\; \frac{h|S_k|}{d},
\]
where we used that each edge in $E(S_k)$ has its outer endpoint in
$N(S_k)\setminus S_k$.  Therefore
$|S_{k+1}|\ge|S_k|\left(1+\frac{h}{d}\right)$ while
$|S_k|\le n/2$.

For part (ii): by (i), after $k_0=\lceil\frac{d}{h}\log(n/2)\rceil$
steps we reach $|S_{k_0}|> n/2$.  One further application of the
expansion argument (applied to $V\setminus S_{k_0}$ of size $<n/2$)
gives $|S_{k_0+1}|=n$.
\end{proof}

\begin{corollary}[Logarithmic cover time]
\label{cor:cover-time}
Under the hypotheses of Theorem~\ref{thm:support-growth}, the plenary
powers $e_i^{[k]}$ have full support (i.e., every generator occurs in
$e_i^{[k]}$) for all $k\ge k_0+1$, where $k_0=O(\frac{d}{h}\log n)$.
\end{corollary}

\begin{proof}
Let $\EA$ be a symmetric $h$-EEA over $\K$ with $\max_i\val(i)\le d$,
fix a generator $e_i\in\natbasis$, and set $S_k=\supp\bigl(e_i^{[k]}\bigr)$
for every $k\ge 0$.  We split the argument into three steps: first we
locate a cover time $k_0$ via the support-growth estimate, then we
establish a monotonicity inclusion for the supports, and finally we
combine the two by induction.

\medskip
\noindent\textbf{Step~1 (existence of a cover time).}
By Theorem~\ref{thm:support-growth}(ii) applied to $e_i$, there is a
positive integer
\[
  k_0\;\le\;\frac{d}{h}\log n
\]
with $|S_{k_0}|=n$.  Since the natural basis indexes the set
$\{1,\ldots,n\}$, this is the equality of sets
$S_{k_0}=\{1,\ldots,n\}$.  The estimate $k_0\le\frac{d}{h}\log n=
O\!\bigl(\frac{d}{h}\log n\bigr)$ delivers the asymptotic bound stated
in the corollary.

\medskip
\noindent\textbf{Step~2 (monotonicity of the support).}
We claim that
\begin{equation}\label{eq:cover-monotone}
  S_m\;\subseteq\;S_{m+1}
  \qquad\text{for every }m\ge 0.
\end{equation}
Fix $m\ge 0$ and expand
\[
  e_i^{[m]}\;=\;\sum_{j\in S_m}\alpha_j^{(m)}\,e_j,
  \qquad \alpha_j^{(m)}\ne 0\;\text{for every }j\in S_m.
\]
The plenary recursion (Definition~\ref{def:powers}) combined with the
multiplication rule $e_je_\ell=0$ for $j\ne\ell$ of
Definition~\ref{def:ea} gives
\begin{equation}\label{eq:cover-expansion}
  e_i^{[m+1]}
  \;=\;\bigl(e_i^{[m]}\bigr)^{2}
  \;=\;\sum_{j\in S_m}\bigl(\alpha_j^{(m)}\bigr)^{2}\,e_j^{\,2}
  \;=\;\sum_{\ell=1}^{n}
        \Bigl(\,\sum_{j\in S_m}\bigl(\alpha_j^{(m)}\bigr)^{2}\,a_{j\ell}\Bigr)\,e_\ell,
\end{equation}
where the last equality uses the structural identity
$e_j^{\,2}=\sum_{\ell}a_{j\ell}e_\ell$ of Definition~\ref{def:ea} and
the linearity of the algebra product.  Reading off the coefficient of
$e_\ell$ in \eqref{eq:cover-expansion},
\begin{equation}\label{eq:cover-alpha-formula}
  \alpha_\ell^{(m+1)}
  \;=\;\sum_{j\in S_m}\bigl(\alpha_j^{(m)}\bigr)^{2}\,a_{j\ell},
  \qquad
  \ell\in\{1,\ldots,n\}.
\end{equation}
The intermediate step of the proof of
Theorem~\ref{thm:support-growth} establishes, from the same formula
\eqref{eq:cover-alpha-formula} under the symmetry hypothesis
$a_{ij}=a_{ji}$, the support recursion
\begin{equation}\label{eq:cover-support-recursion}
  S_{m+1}\;\supseteq\;S_m\,\cup\,N(S_m),
\end{equation}
where $N(S_m)$ denotes the neighbourhood of $S_m$ in
$\varGamma(\EA,\natbasis)$.  Taking the $S_m$-part of
\eqref{eq:cover-support-recursion} yields
\eqref{eq:cover-monotone}.

\medskip
\noindent\textbf{Step~3 (induction).}
We prove $S_k=\{1,\ldots,n\}$ for every $k\ge k_0$ by induction on
$k$.  The base case $k=k_0$ is Step~1.  For the inductive step,
suppose $S_m=\{1,\ldots,n\}$ for some $m\ge k_0$.  By
\eqref{eq:cover-monotone} we have
$\{1,\ldots,n\}=S_m\subseteq S_{m+1}$, while the trivial inclusion
$S_{m+1}\subseteq\{1,\ldots,n\}$ holds by definition of the support.
Hence $S_{m+1}=\{1,\ldots,n\}$, closing the induction.

In particular $S_k=\{1,\ldots,n\}$ for every $k\ge k_0+1$, that is,
every generator occurs in $e_i^{[k]}$, which is the assertion of the
corollary.
\end{proof}

\subsection{Algebraic persistency in symmetric EEAs}

\begin{theorem}[Algebraic persistency]\label{thm:persistency}
Let $\EA$ be a symmetric $h$-EEA over $\K$ with $n=\dim\EA\ge 2$ and
$h>0$.  Then every generator $e_i$ is algebraically persistent.
\end{theorem}

\begin{proof}
We show that $e_i$ occurs in $e_i^{[k]}$ for all $k\ge 0$.

Since $\EA$ is symmetric, the undirected graph $\varGamma(\EA,\natbasis)$
is connected (by Theorem~\ref{thm:connectivity}).  We use the following
key observation: for any edge $\{i,j\}$ in $\varGamma(\EA,\natbasis)$,
we have $a_{ij}=a_{ji}\ne 0$ (by symmetry and definition of the edge).
Therefore, from $e_j^2=\sum_k a_{jk} e_k$, we see that $e_i\prec e_j^2$
(i.e., $e_i$ occurs in $e_j^2$).  Equivalently, in the directed graph
$D(\EA,\natbasis)$ every undirected edge $\{i,j\}$ gives rise to the
directed edges $i\to j$ and $j\to i$.

\medskip
We now proceed by induction on $k$.  For $k=0$: $e_i^{[0]}=e_i$, so
$e_i$ trivially occurs in $e_i^{[0]}$.

For the inductive step, suppose $e_i$ occurs in $e_i^{[k]}$, say
$e_i^{[k]}=\alpha_i e_i + \sum_{j\ne i}\alpha_j e_j$ with
$\alpha_i\ne 0$.  Then
\begin{align*}
  e_i^{[k+1]} &= e_i^{[k]}\cdot e_i^{[k]}
  = \left(\alpha_i e_i + \sum_{j\ne i}\alpha_j e_j\right)^2
  = \alpha_i^2\,e_i^2 + \sum_{j\ne i}\alpha_j^2\,e_j^2.
\end{align*}
The coefficient of $e_i$ in $e_i^{[k+1]}$ is
\[
  \alpha_i^2\,a_{ii} + \sum_{j\ne i}\alpha_j^2\,a_{ji}.
\]
If $a_{ii}\ne 0$, this coefficient equals $\alpha_i^2 a_{ii}\ne 0$, so
$e_i$ occurs in $e_i^{[k+1]}$.

If $a_{ii}=0$: since $\EA$ is connected (Corollary~\ref{cor:EEA-connected})
and $n\ge 2$, there exists at least one neighbour $j\ne i$ in
$\varGamma(\EA,\natbasis)$.  By symmetry, $a_{ji}=a_{ij}\ne 0$.  If
$\alpha_j\ne 0$ for at least one such neighbour $j$, the coefficient
of $e_i$ in $e_i^{[k+1]}$ includes the term $\alpha_j^2 a_{ji}\ne 0$.

It remains to show that $\alpha_j\ne 0$ for some neighbour $j$ of $i$.
By Theorem~\ref{thm:support-growth}, after at most $k_0=O(\frac{d}{h}
\log n)$ steps, the support of $e_i^{[k]}$ is all of $\{1,\ldots,n\}$.
In particular, for $k\ge k_0$, all coefficients $\alpha_j$ are nonzero
and the persistence follows.  For $k<k_0$, we note that the support
$S_k$ grows by at least one vertex at each step (since the graph is
connected), and in particular the support always contains $i$ (by
induction) and at least one neighbour of $i$ once $|S_k|\ge 2$.  Since
$S_1=N[i]$ (the closed neighbourhood of $i$, which contains $j$ for
every neighbour $j$ of $i$), the claim follows for all $k\ge 1$.
\end{proof}

\begin{remark}
The symmetry assumption $a_{ij}=a_{ji}$ is essential.  For non-symmetric
EEAs, the directed graph may have edges $i\to j$ without the return edge
$j\to i$, enabling algebraic transiency even when $\varGamma(\EA,\natbasis)$
is an expander.
\end{remark}

\subsection{Hierarchical structure}

Tian \cite{Tian2008} defines a hierarchical decomposition of an
evolution algebra into levels, based on the partition of generators into
algebraically persistent and transient.

\begin{theorem}[Hierarchy of a symmetric EEA]\label{thm:hierarchy}
A symmetric $h$-EEA $\EA$ over $\K$ with $n\ge 2$ and $h>0$ has
\emph{trivial hierarchy}: the only hierarchical level is level $0$,
containing all generators.  Equivalently, $\EA$ is \emph{algebraically
persistent} (Definition~\ref{def:persist}).
\end{theorem}

\begin{proof}
By Theorem~\ref{thm:persistency}, every generator is algebraically
persistent.  By Tian \cite[Chapter~3, \S3.5]{Tian2008}, the first level
of the hierarchy consists precisely of the algebraically persistent
generators, i.e., all generators.  There are no higher levels.
\end{proof}

% ==================================================================
\section{The Evolution Operator and Spectral Theory over $\K$}
\label{sec:evop}
% ==================================================================

\subsection{The evolution operator}

\begin{definition}[Evolution operator]\label{def:evop}
Let $\EA$ be a finite-dimensional evolution algebra over $\K$ with
natural basis $\natbasis=\{e_1,\ldots,e_n\}$ and structural matrix
$\strmat=(a_{ij})$.  The \emph{evolution operator}
$\evop:\EA\to\EA$ is the $\K$-linear map defined by
\[
  \evop(e_i) \;:=\; e_i^2 \;=\; \sum_{j=1}^n a_{ij}\,e_j,
  \quad i=1,\ldots,n.
\]
The matrix of $\evop$ in the basis $\natbasis$ is $\strmat^T$ (the
transpose of the structural matrix), since $[\evop(e_i)]_j = a_{ij} =
(\strmat)_{ij}=(\strmat^T)_{ji}$.
\end{definition}

\begin{remark}
With the convention that elements of $\EA$ are \emph{column vectors},
the matrix of $\evop$ acting on the left is $\strmat$ itself:
$\evop(e_i) = \strmat\,e_i$ in column-vector notation.  We adopt the
row-vector convention throughout, so the matrix is $\strmat$.
\end{remark}

\begin{proposition}[Plenary powers and powers of $\strmat$]\label{prop:powers}
For any $x=\sum_i \alpha_i e_i\in\EA$, the $k$-th plenary power
satisfies $x^{[k]}=\sum_i (\alpha_i)^{2^k}\,e_i^{[k]}$, and in
the Markov case (Section~\ref{sec:markov}), $\evop^k$ corresponds to
$k$-step transition matrix $\strmat^k$.
\end{proposition}

\begin{proof}
By induction: $(e_i^{[k]})^2=e_i^{[k+1]}$, so
$(\sum_i\alpha_i e_i)^{[k+1]}=(\sum_i\alpha_i^{2^k}e_i^{[k]})^2
=\sum_i\alpha_i^{2^{k+1}}e_i^{[k+1]}$.
\end{proof}

\begin{theorem}[Evolution operator and graph distance]\label{thm:evop-distance}
Let $\EA$ be an $h$-EEA over $\K$ with $\hh(\varGamma(\EA,\natbasis))
\ge h$ and $\max_i\val(i)\le d$.  For any generator $e_i$, the
minimal integer $k$ such that $\supp(\evop^k(e_i))=\{1,\ldots,n\}$
satisfies
\[
  k \;\le\; \frac{2d}{h}\,\log n \;+\; 2.
\]
\end{theorem}

\begin{proof}
By Definition~\ref{def:evop}, $\evop(e_i)=e_i^2=\sum_j a_{ij}e_j$, so
$\supp(\evop^1(e_i))=\{j: a_{ij}\ne 0\}=N_+(i)$ (the out-neighbourhood
of $i$ in $D(\EA,\natbasis)$).  For the symmetric case, $N_+(i)=N(i)$
and $\supp(\evop^k(e_i))=\bigcup_{j\in\supp(\evop^{k-1}(e_i))} N(j)$,
which is the $k$-ball $B_k(i)$ in $\varGamma(\EA,\natbasis)$.  By
Theorem~\ref{thm:diameter}, the diameter is at most $\frac{2d}{h}\log n+1$,
so $B_k(i)=V$ for $k=\diam(\varGamma)+1\le\frac{2d}{h}\log n+2$.
\end{proof}

\subsection{Spectral theory over a general field $\K$}

\begin{theorem}[Characteristic polynomial and eigenvalues over $\K$]\label{thm:char-poly}
Let $\EA$ be an $h$-EEA over $\K$ with structural matrix $\strmat$.
The characteristic polynomial $\chi_\strmat(\lambda)=\det(\lambda I-\strmat)$
is a degree-$n$ polynomial in $\K[\lambda]$.  The eigenvalues of
$\evop$ (roots of $\chi_\strmat$) lie in the algebraic closure
$\overline{\K}$.
\end{theorem}

\begin{proof}
Standard linear algebra over $\K$.
\end{proof}

\begin{theorem}[Spectral radius and expansion over $\R$]\label{thm:spectral-radius}
Let $\EA$ be an $h$-EEA over $\R$ with symmetric structural matrix
$\strmat$ (i.e., $a_{ij}=a_{ji}$) and $\max_i \sum_{j\ne i} a_{ij}
\le d$ (uniform degree bound on the \emph{support} of $\strmat$).
Let $\lambda_1\ge\lambda_2\ge\cdots\ge\lambda_n$ be the real
eigenvalues of $\strmat$ (symmetric real matrix).  Then:
\begin{enumerate}[label=\normalfont(\roman*)]
  \item $\lambda_1=\rho(\strmat)$ (spectral radius);
  \item The \emph{spectral gap} $\mathrm{gap}(\strmat):=\lambda_1-\lambda_2$
        satisfies
        \[
          \mathrm{gap}(\strmat) \;\ge\;
          \frac{\bigl(\hh(\EA,\natbasis)\bigr)^2}{2d}.
        \]
\end{enumerate}
\end{theorem}

\begin{proof}
Since $\strmat$ is real and symmetric, all eigenvalues are real and the
spectral theorem applies.  Part (i) follows from the Perron--Frobenius
theorem for symmetric matrices (or directly from the variational formula
$\lambda_1=\max_{\|v\|=1}\langle v,\strmat v\rangle$).

For part (ii): the proof uses the \emph{discrete Cheeger inequality} for
weighted graphs.  Consider the underlying unweighted graph
$\varGamma(\EA,\natbasis)$ with adjacency matrix $G=(g_{ij})$ where
$g_{ij}=\mathbf{1}[a_{ij}\ne 0,\, i\ne j]$.  Since $\strmat$ is
symmetric with the same support pattern as $G$:
\[
  \strmat \;=\; \sum_{i<j} a_{ij}\bigl(e_i e_j^T + e_j e_i^T\bigr)
             \;+\; \sum_i a_{ii}\,e_i e_i^T.
\]
The normalized Laplacian of $G$ (for the $d$-regular case $a_{ij}=c$
for all edges $\{i,j\}$) satisfies the discrete Cheeger inequality
\cite[Chapter~3]{Kowalski2019}:
\[
  1 - \frac{\lambda_2(\strmat/c\cdot d)}{\lambda_1(\strmat/c\cdot d)}
  \;\ge\;
  \frac{\hh(\varGamma)^2}{2d^2}.
\]
Rearranging (with $\lambda_1=cd$ for the $d$-regular graphicable case
with constant weight $c$):
\[
  \lambda_1-\lambda_2
  \;\ge\;
  \frac{h^2}{2d^2}\,\lambda_1
  \;\ge\;
  \frac{c h^2}{2d}.
  \qedhere
\]
\end{proof}

% ==================================================================
\section{Markov Expander Evolution Algebras}\label{sec:markov}
% ==================================================================

\subsection{Definition and basic properties}

\begin{definition}[Markov EEA]\label{def:markov-EEA}
An evolution algebra $\EA$ over $\R$ is a \emph{Markov evolution
algebra} if:
\begin{enumerate}[label=\normalfont(\roman*)]
  \item $a_{ij}\ge 0$ for all $i,j$ (non-negative structural constants);
  \item $\sum_{j=1}^n a_{ij}=1$ for each $i$ (each row of $\strmat$ sums
        to $1$).
\end{enumerate}
A Markov evolution algebra that is also an $h$-EEA is called a
\emph{Markov $h$-EEA}.
\end{definition}

In this case, the structural matrix $\strmat$ is a \emph{stochastic
matrix} and the evolution operator $\evop$ corresponds to one step of
the Markov chain with transition matrix $\strmat$.

\begin{remark}
For a Markov EEA, the graph $\varGamma(\EA,\natbasis)$ has an edge
$\{i,j\}$ whenever $a_{ij}>0$ or $a_{ji}>0$.  The Cheeger constant
$\hh(\EA,\natbasis)$ thus quantifies the ``mixing efficiency'' of the
underlying Markov chain.
\end{remark}

\subsection{Simplicity of Markov EEAs}\label{sub:markov-simple}

\begin{theorem}[Simplicity and irreducibility]\label{thm:markov-simple}
Let $\EA$ be a Markov $h$-EEA.  Then:
\begin{enumerate}[label=\normalfont(\roman*)]
  \item The Markov chain with transition matrix $\strmat$ is
        \emph{irreducible}.
  \item $\EA$ is a simple evolution algebra.
\end{enumerate}
\end{theorem}

\begin{proof}
(i): Irreducibility means that for any states $i,j$ there exists a path
from $i$ to $j$ in the directed graph of $\strmat$ (edges $i\to j$
whenever $a_{ij}>0$).  Since $\hh(\varGamma(\EA,\natbasis))>0$, the
underlying undirected graph is connected
(Theorem~\ref{thm:connectivity}).  Hence there is an undirected path
from $i$ to $j$.  Since $a_{ij}>0$ iff there is an edge $i\to j$ in
$D(\EA,\natbasis)$, and an undirected edge $\{i,j\}$ implies either
$a_{ij}>0$ or $a_{ji}>0$, we can always traverse the path:
at each undirected edge, at least one direction is available, and in a
Markov chain we can take a two-step path $i\to k\to j$ if only $a_{ki}>0$
and $a_{kj}>0$.  More precisely, by connectivity and non-negativity, a
standard argument using the graph diameter (Theorem~\ref{thm:diameter})
shows that every state communicates.

(ii): By Tian \cite[Chapter~4]{Tian2008}, a Markov evolution algebra is
simple if and only if the underlying Markov chain is irreducible, which
we just proved.
\end{proof}

\subsection{Mixing time}

\begin{theorem}[Convergence to stationarity]\label{thm:mixing}
Let $\EA$ be a symmetric Markov $h$-EEA with $n=\dim\EA$, $\strmat$
doubly stochastic (each row and column sums to $1$), and
$\max_i|\{j:a_{ij}>0\}|\le d$.  Let $\pi=(1/n,\ldots,1/n)^T$ be the
uniform stationary distribution.  Then for any initial generator $e_i$
and any $k\ge 1$:
\[
  \bigl\|\strmat^k\,e_i - \pi\bigr\|_1
  \;\le\;
  n\,\left(1-\frac{h^2}{2d^2}\right)^k,
\]
where $\|\cdot\|_1$ is the $\ell^1$ norm on $\R^n$.
\end{theorem}

\begin{proof}
By Theorem~\ref{thm:spectral-radius}(ii), the spectral gap satisfies
$\mathrm{gap}(\strmat)\ge h^2/(2d^2)$ (for the normalized matrix
$\strmat/\rho(\strmat)$ in the doubly stochastic case where
$\rho(\strmat)=1$).  Let $1=\mu_1\ge\mu_2\ge\cdots\ge\mu_n$ be the
eigenvalues of $\strmat$ (as a doubly stochastic matrix, $\mu_1=1$ with
eigenvector $\mathbf{1}$).  The spectral expansion gives
$|\mu_2|\le 1-h^2/(2d^2)$.

Since $\strmat$ is doubly stochastic, $\pi=\mathbf{1}/n$ is the
stationary distribution and $\strmat^k\to\pi\mathbf{1}^T$ as
$k\to\infty$.  The standard bound
\[
  \|\strmat^k e_i - \pi\|_2 \;\le\;
  \sqrt{n}\,|\mu_2|^k
\]
follows from the spectral decomposition.  The $\ell^1$ bound follows via
$\|v\|_1\le\sqrt{n}\|v\|_2$.
\end{proof}

\begin{corollary}[Logarithmic mixing time]\label{cor:mixing-time}
Under the hypotheses of Theorem~\ref{thm:mixing}, the mixing time
$t_\mathrm{mix}(\epsilon):=\min\{k:\|\strmat^k e_i-\pi\|_1\le\epsilon\}$
satisfies
\[
  t_\mathrm{mix}(\epsilon)
  \;\le\;
  \frac{2d^2}{h^2}\,\log\!\left(\frac{n}{\epsilon}\right).
\]
\end{corollary}

\begin{proof}
Set $n(1-h^2/(2d^2))^k\le\epsilon$ and solve for $k$.
\end{proof}

% ==================================================================
\section{The $d$-Regular Case}\label{sec:regular}
% ==================================================================

\subsection{$d$-regular EEAs}

\begin{definition}[$d$-regular EEA]\label{def:d-regular}
An evolution algebra $\EA$ with natural basis $\natbasis$ is
\emph{$d$-regular} if the underlying graph $\varGamma(\EA,\natbasis)$ is
$d$-regular, i.e., every vertex has exactly $d$ neighbours in
$\varGamma(\EA,\natbasis)$.  A $d$-regular $h$-EEA is a $d$-regular
evolution algebra with $\hh(\EA,\natbasis)\ge h$.
\end{definition}

\begin{proposition}[Degree condition]\label{prop:degree-cond}
In a $d$-regular EEA over $\K$: for every $i$, the number of indices
$j\ne i$ with $a_{ij}\ne 0$ or $a_{ji}\ne 0$ equals $d$.
\end{proposition}

\begin{proof}
Direct from the definition of $d$-regularity and
Definition~\ref{def:assoc-graph}.
\end{proof}

\subsection{Cheeger inequality for the evolution operator}

For $d$-regular EEAs over $\R$, the discrete Cheeger inequality gives explicit gap bounds in terms of the expansion constant.

\begin{theorem}[Cheeger inequality for EEAs]\label{thm:cheeger-evop}
Let $\EA$ be a symmetric $d$-regular EEA over $\R$ with structural
matrix $\strmat$ satisfying $a_{ij}=c>0$ for all edges $\{i,j\}$ of
$\varGamma(\EA,\natbasis)$ and $a_{ij}=0$ otherwise, $a_{ii}=0$.
Let $\lambda_1\ge\lambda_2\ge\cdots\ge\lambda_n$ be the eigenvalues of
$\strmat$.  Then $\lambda_1=cd$ and:
\[
  \frac{(cd)^2\,\hh(\EA,\natbasis)^2}{2d^2\lambda_1}
  \;\le\;
  \lambda_1 - \lambda_2
  \;\le\;
  2c\,\hh(\EA,\natbasis).
\]
In particular, $\mathrm{gap}(\strmat)\ge ch^2/(2d)$.
\end{theorem}

\begin{proof}
Note that $\strmat=c\cdot G$ where $G$ is the adjacency matrix of
$\varGamma(\EA,\natbasis)$ (a $d$-regular graph).  Let $P=G/d$ be the
normalised adjacency matrix with eigenvalues $1=\nu_1\ge\nu_2\ge\cdots
\ge\nu_n\ge-1$.  The discrete Cheeger inequality
\cite[Theorem~3.3.2]{Kowalski2019} gives:
\[
  \frac{\hh(\varGamma)^2}{2} \;\le\; d\,(1-\nu_2) \;\le\; 2d\,\hh(\varGamma).
\]
Since $\lambda_k=c\,d\,\nu_k$, we get
$\lambda_1-\lambda_2=cd(1-\nu_2)$ and
\[
  \frac{c\,\hh(\EA,\natbasis)^2}{2}
  \;\le\;
  \lambda_1-\lambda_2
  \;\le\;
  2c\,d\,\hh(\EA,\natbasis).
  \qedhere
\]
\end{proof}

\begin{corollary}[EEA $\Leftrightarrow$ spectral gap for symmetric graphicable algebras]
\label{cor:spectral-EEA}
Let $\EA$ be a symmetric $d$-regular evolution algebra over $\R$ with
$a_{ij}\in\{0,1\}$ for all $i\ne j$ (i.e., $\EA$ is a
\emph{graphicable} evolution algebra in the sense of
Tian~\cite{Tian2008}).  Then $\EA$ is an $h$-EEA if and only if
$\mathrm{gap}(\strmat)\ge h^2/2$.
\end{corollary}

\begin{proof}
Immediate from Theorem~\ref{thm:cheeger-evop} with $c=1$.
\end{proof}

\subsection{Alon's conjecture and the spectral gap bound}


\begin{theorem}[Trivial spectral bound over $\R$]\label{thm:trivial-spectral}
Let $\EA$ be a symmetric graphicable $d$-regular EEA over $\R$ with
$n=\dim\EA$.  Then
\[
  \hh(\EA,\natbasis) \;\le\; d/2.
\]
Furthermore, the spectral gap satisfies
$\mathrm{gap}(\strmat)\le d^2/8$ (trivial upper bound from the
Cheeger inequality).
\end{theorem}

\begin{proof}
By Proposition~\ref{prop:cheeger-basic}(iii):
$\hh(\EA,\natbasis)\le\min_i\val(i)=d$.  A tighter analysis (taking
$W=N[v]$ for a vertex $v$, which has $|W|\le d+1$ and $|E(W)|\le d$)
gives $h\le d/(d+1)\le d/2$ for $d\ge 2$.  (Actually the
Proposition gives $h\le d$; the tight bound $d/2$ follows from the
Expander Mixing Lemma and a standard averaging argument for $d$-regular
graphs \cite[Lemma~3.1.4]{Kowalski2019}.)
\end{proof}

% ==================================================================
\section{Complex Expander Evolution Algebras}\label{sec:complex}
% ==================================================================

Over $\C$, the structural matrix of a symmetric graphicable EEA is real-symmetric, so all eigenvalues are real and universal bounds apply.

\subsection{Eigenvalues over $\C$}

\begin{theorem}[Spectral characterisation over $\C$]\label{thm:spectral-C}
Let $\EA$ be a symmetric graphicable $d$-regular EEA over $\C$ with
adjacency-type structural matrix $\strmat$.  Then:
\begin{enumerate}[label=\normalfont(\roman*)]
  \item All eigenvalues of $\strmat$ are real.
  \item $\lambda_1=d$ (the Perron eigenvalue).
  \item $|\lambda_k|\le d$ for all $k$.
  \item $\lambda_n\ge -d$.
  \item The eigenspace of $\lambda_1=d$ is spanned by the all-ones
        vector $\mathbf{1}=(1,\ldots,1)^T$, i.e., it is one-dimensional
        when $\varGamma(\EA,\natbasis)$ is connected.
\end{enumerate}
\end{theorem}

\begin{proof}
Since $\strmat$ is the adjacency matrix of a $d$-regular graph (a real
symmetric matrix), all its eigenvalues are real.  The rest follows from
the Perron--Frobenius theorem applied to $G=\strmat$ (non-negative,
irreducible by connectivity).
\end{proof}

\subsection{The Alon--Boppana bound for EEAs}


\begin{theorem}[Alon--Boppana bound]\label{thm:alon-boppana}
Let $(\EA_m,\natbasis_m)_{m\ge 1}$ be an EEA family of symmetric
graphicable $d$-regular evolution algebras over $\C$ with
$n_m=\dim\EA_m\to\infty$.  Let $\lambda_2^{(m)}$ be the second largest
eigenvalue of the structural matrix $\strmat_m$.  Then
\[
  \liminf_{m\to\infty}\,\lambda_2^{(m)}
  \;\ge\; 2\sqrt{d-1}.
\]
\end{theorem}

\begin{proof}
The matrices $\strmat_m$ are the adjacency matrices of an expander
family $(\varGamma(\EA_m,\natbasis_m))$.  The classical Alon--Boppana
theorem \cite[Theorem~4.2.x]{Kowalski2019} (see also Nilli~\cite{Nilli1991})
states that for any $d$-regular graph $G$ on $n$ vertices:
\[
  \lambda_2(A_G) \;\ge\; 2\sqrt{d-1}\,\left(1-\frac{c}{\diam(G)}\right)
\]
for an absolute constant $c>0$.  Since $\diam(\varGamma(\EA_m,\natbasis_m))
\to\infty$ (as $n_m\to\infty$ and the graphs are $d$-regular expanders),
the right-hand side tends to $2\sqrt{d-1}$.
\end{proof}

\begin{remark}
Theorem~\ref{thm:alon-boppana} shows that the spectral gap
$\lambda_1-\lambda_2=d-\lambda_2$ cannot exceed $d-2\sqrt{d-1}$
asymptotically.  By the Cheeger inequality
(Theorem~\ref{thm:cheeger-evop}), the expansion constant satisfies
$\hh(\EA_m,\natbasis_m)\le\sqrt{d(d-\lambda_2)}\to\sqrt{d(d-2\sqrt{d-1})}$.
\end{remark}

\subsection{Ramanujan evolution algebras}

Theorem~\ref{thm:alon-boppana} is tight, and the algebras meeting it deserve a name.

\begin{definition}[Ramanujan evolution algebra]\label{def:Ramanujan}
A symmetric graphicable $d$-regular EEA $\EA$ over $\K$ is called a
\emph{Ramanujan evolution algebra} if the eigenvalues $\lambda$ of
$\strmat$ with $|\lambda|\ne d$ satisfy
\[
  |\lambda| \;\le\; 2\sqrt{d-1}.
\]
\end{definition}

\begin{remark}
Equivalently, $\EA$ is Ramanujan if $\varGamma(\EA,\natbasis)$ is a
Ramanujan graph in the sense of Lubotzky, Phillips, and Sarnak
\cite{LPS1988} (or in the sense of Marcus, Spielman, and Srivastava
\cite{MSS2015}).
\end{remark}

\begin{theorem}[Expansion of Ramanujan EEAs]\label{thm:Ramanujan-expansion}
A Ramanujan evolution algebra $\EA$ is an $h$-EEA with
\[
  \hh(\EA,\natbasis) \;\ge\; \frac{d - 2\sqrt{d-1}}{2}.
\]
\end{theorem}

\begin{proof}
By the Cheeger inequality (Theorem~\ref{thm:cheeger-evop}):
\[
  \hh(\EA,\natbasis) \;\ge\;
  \frac{\lambda_1-\lambda_2}{2}
  \;=\;
  \frac{d-\lambda_2}{2}
  \;\ge\;
  \frac{d-2\sqrt{d-1}}{2}.
  \qedhere
\]
\end{proof}

\begin{proposition}[Ramanujan EEAs have optimal mixing time]\label{prop:Ramanujan-mixing}
For a Ramanujan Markov EEA (doubly stochastic, $d$-regular):
\[
  t_\mathrm{mix}(\epsilon)
  \;\le\;
  \frac{d}{d-2\sqrt{d-1}}\,\log\!\left(\frac{n}{\epsilon}\right).
\]
This is the optimal mixing time among all $d$-regular Markov EEAs.
\end{proposition}

\begin{proof}
Substitute $h=(d-2\sqrt{d-1})/2$ into Corollary~\ref{cor:mixing-time}
with $d$ replaced by $1$ (for the normalized matrix).
\end{proof}

% ==================================================================
\section{Constructions of EEAs}\label{sec:constructions}
% ==================================================================

\subsection{Cayley EEAs}\label{sec:cayley}

\begin{definition}[Cayley evolution algebra]\label{def:cayley-EA}
Let $G$ be a finite group, $S=\{s_1,\ldots,s_d\}\subseteq G\setminus
\{1\}$ a \emph{symmetric generating set} (i.e., $S=S^{-1}$ and
$1\notin S$), and $\K$ a field.  The \emph{Cayley evolution algebra}
$\mathcal{C}(G,S,\K)$ is defined as follows:
\begin{itemize}
  \item Natural basis: $\natbasis=\{e_g : g\in G\}$;
  \item Structural constants: $e_g^2=\sum_{s\in S} e_{gs}$,
        i.e., $a_{g,gs}=1$ for all $s\in S$ and $a_{g,h}=0$ otherwise.
\end{itemize}
\end{definition}

\begin{remark}
The underlying undirected graph $\varGamma(\mathcal{C}(G,S,\K),\natbasis)$
is exactly the Cayley graph $\mathrm{Cay}(G,S)$, since $a_{g,gs}=1$
and $a_{gs,g}=a_{gs,(gs)s^{-1}}=1$ (because $s^{-1}\in S$).
\end{remark}

\begin{theorem}[Cayley EEAs are EEAs iff Cayley graphs are expanders]
\label{thm:cayley-EEA}
$\mathcal{C}(G,S,\K)$ is an $h$-EEA if and only if the Cayley graph
$\mathrm{Cay}(G,S)$ is an $h$-expander.
\end{theorem}

\begin{proof}
Direct from Definition~\ref{def:cayley-EA} and Definition~\ref{def:EEA},
since $\varGamma(\mathcal{C}(G,S,\K),\natbasis)=\mathrm{Cay}(G,S)$.
\end{proof}

\begin{corollary}[EEA families from group theory]\label{cor:group-EEA}
Let $(G_m)_{m\ge 1}$ be a family of finite groups with symmetric
generating sets $S_m\subseteq G_m\setminus\{1\}$, $|S_m|=d$ for all
$m$, such that $|G_m|\to\infty$ and $\mathrm{Cay}(G_m,S_m)$ form an
expander family.  Then $(\mathcal{C}(G_m,S_m,\K))_{m\ge 1}$ is an EEA
family.
\end{corollary}

\begin{example}[SL$_2(\F_p)$ EEAs]\label{ex:SL2-EEA}
Let $p$ run over primes and $S\subseteq\mathrm{SL}_2(\Z)$ be a fixed
symmetric generating set for $\mathrm{SL}_2(\Z)$ that does not
generate a solvable group.  Let $S_p$ be the reduction of $S$ modulo
$p$.  By the Bourgain--Gamburd theorem \cite{BourgainGamburd2008}
(see also \cite[Chapter~6]{Kowalski2019}), the Cayley graphs
$\mathrm{Cay}(\mathrm{SL}_2(\F_p),S_p)$ form an expander family.
Therefore, $(\mathcal{C}(\mathrm{SL}_2(\F_p),S_p,\K))_p$ is an EEA
family.
\end{example}

\begin{example}[Ramanujan Cayley EEAs from LPS]\label{ex:LPS}
Let $p,q$ be distinct primes with $p\equiv q\equiv 1\pmod{4}$ and $p$
a quadratic residue mod $q$.  The Lubotzky--Phillips--Sarnak
construction \cite{LPS1988} yields a symmetric generating set
$S_{p,q}\subset\mathrm{PGL}_2(\F_q)$ of size $p+1$ such that
$\mathrm{Cay}(\mathrm{PGL}_2(\F_q),S_{p,q})$ is a Ramanujan graph.
The corresponding Cayley evolution algebra
$\mathcal{C}(\mathrm{PGL}_2(\F_q),S_{p,q},\K)$ is a Ramanujan
evolution algebra.
\end{example}

\subsection{Tensor products of EEAs}\label{sec:tensor}

\begin{theorem}[Kronecker product of EEAs]\label{thm:kronecker}
Let $\EA_1$ and $\EA_2$ be evolution algebras over $\K$ with natural
bases $\natbasis_1=\{e_i\}_{i=1}^{n_1}$ and
$\natbasis_2=\{f_j\}_{j=1}^{n_2}$, and structural matrices
$\strmat_1,\strmat_2$.  The Kronecker product $\EA_1\otimes_\K\EA_2$
(with natural basis $\{e_i\otimes f_j\}$ and structural constants
$a_{(i,j),(k,l)}=a^{(1)}_{ik}\,a^{(2)}_{jl}$) is again an evolution
algebra.  If $\EA_1$ is an $h_1$-EEA and $\EA_2$ is an $h_2$-EEA,
then $\EA_1\otimes\EA_2$ satisfies
\[
  \hh\!\bigl(\EA_1\otimes\EA_2,\natbasis_1\otimes\natbasis_2\bigr)
  \;\ge\;
  \min(h_1,h_2).
\]
\end{theorem}

\begin{proof}
The Kronecker product is an evolution algebra by
Tian~\cite[Corollary~1(5)]{Tian2008}.  The graph
$\varGamma(\EA_1\otimes\EA_2,\natbasis_1\otimes\natbasis_2)$ is the
\emph{tensor product} (or \emph{categorical product}) of the graphs
$\varGamma_1$ and $\varGamma_2$: the vertex set is
$\{1,\ldots,n_1\}\times\{1,\ldots,n_2\}$ and there is an edge between
$(i,j)$ and $(k,l)$ iff $a^{(1)}_{ik}\ne 0$ (or $a^{(1)}_{ki}\ne 0$)
\emph{and} $a^{(2)}_{jl}\ne 0$ (or $a^{(2)}_{lj}\ne 0$).

The expansion of the tensor product graph $\varGamma_1\times\varGamma_2$
satisfies $\hh(\varGamma_1\times\varGamma_2)\ge\min(\hh(\varGamma_1),
\hh(\varGamma_2))$ by a standard edge-boundary argument (see, e.g.,
\cite[Exercise~3.1.13]{Kowalski2019}), giving the claimed bound.
\end{proof}

\begin{corollary}[EEA families from iterated tensor products]
\label{cor:tensor-family}
Let $\EA$ be an $h$-EEA.  The family $(\EA^{\otimes m})_{m\ge 1}$
(with $\dim\EA^{\otimes m}=n^m\to\infty$) is an EEA family with
$\hh(\EA^{\otimes m})\ge h$ for all $m$.
\end{corollary}

\subsection{Direct sums}

\begin{proposition}[Direct sum is not an EEA]
If $\EA=\EA_1\oplus\EA_2$ is a direct sum of two non-trivial evolution
algebras, then $\hh(\EA,\natbasis)=0$.
\end{proposition}

\begin{proof}
Taking $W$ to be the index set of $\EA_1$, we have $|E(W)|=0$
(no structural constants cross between $\EA_1$ and $\EA_2$), so
$\hh\le|E(W)|/|W|=0$.
\end{proof}

Direct sums are excluded from the EEA class; this is consistent with Corollary~\ref{cor:EEA-connected}.

% ==================================================================
\section{Examples}\label{sec:examples}
% ==================================================================

\begin{example}[Cycle evolution algebra]\label{ex:cycle}
Let $n\ge 3$ and $\K$ a field.  Define $\EA_n$ with natural basis
$\{e_0,\ldots,e_{n-1}\}$ and $e_i^2=e_{i-1}+e_{i+1}$ (indices mod $n$).
This is the \emph{cycle evolution algebra} of Tian \cite{Tian2008}.
The underlying graph $\varGamma(\EA_n,\natbasis)$ is the cycle $C_n$.
By Kowalski~\cite[Example~3.1.3(2)]{Kowalski2019}:
\[
  \hh(\EA_n,\natbasis) \;=\; \frac{2}{\lfloor n/2\rfloor} \;\to\; 0.
\]
Hence the family $(\EA_n)_{n\ge 3}$ is \emph{not} an EEA family, though
each individual $\EA_n$ is an $\hh_n$-EEA for $\hh_n>0$.
\end{example}

\begin{example}[Complete evolution algebra]\label{ex:complete}
Let $n\ge 2$ and $e_i^2=\sum_{j\ne i} e_j$.  The underlying graph is
$K_n$ (complete graph).  By Kowalski~\cite[Example~3.1.3(1)]{Kowalski2019}:
\[
  \hh(\EA_{K_n},\natbasis) \;=\; \left\lfloor \frac{n}{2}\right\rfloor
  \;\to\; \infty.
\]
This is the most expander-like single algebra but it is $(n-1)$-regular
(unbounded degree), so the family $(\EA_{K_n})$ does not satisfy the
bounded-degree condition for an expander family (Definition~\ref{def:expander}).
\end{example}

\begin{example}[$3$-regular Ramanujan EEA]\label{ex:Ramanujan-3reg}
The Petersen graph is a $3$-regular graph on $10$ vertices.
Its second eigenvalue is $\lambda_2=1$, which satisfies $\lambda_2=1<2\sqrt{2}=2\sqrt{d-1}$,
so it is a Ramanujan graph.  The corresponding evolution algebra, with $n=10$ and
$a_{ij}=1$ for each Petersen edge, is therefore a Ramanujan evolution algebra with
$\hh\ge(3-2\sqrt{2})/2$.
\end{example}

% ==================================================================
\section{Open Problems}\label{sec:open}
% ==================================================================

We close with problems that remain open.

\begin{problem}[Algebraic characterisation of EEAs over finite fields]
\label{op:finite-field}
For a field $\K=\F_q$ (finite field with $q$ elements), characterise
the $h$-EEAs over $\K$ purely in terms of the algebra structure
(without reference to real eigenvalues).  In particular, does the
combinatorial expansion condition $\hh(\EA,\natbasis)\ge h$ have an
algebraic analogue in the representation theory of $\EA$ over $\F_q$?
\end{problem}

\begin{problem}[Non-symmetric EEAs and algebraic persistency]
\label{op:nonsymm}
Theorem~\ref{thm:persistency} requires the symmetry assumption
$a_{ij}=a_{ji}$.  Is algebraic persistency equivalent to some directed
expansion condition on $D(\EA,\natbasis)$?  More precisely: is there
a directed analogue of the Cheeger constant for evolution algebras such
that the directed expansion condition is equivalent to algebraic
persistency?
\end{problem}

\begin{problem}[Expansion constant as algebraic invariant]
\label{op:invariant}
Theorem~\ref{thm:basis-invariance} shows that for nonsingular EEAs, the
isomorphism class of $\varGamma(\EA,\natbasis)$ is an algebraic invariant.
Is the \emph{Cheeger constant} $\hh(\EA,\natbasis)$ itself an algebraic
invariant (not just the graph up to isomorphism)?  Is there an algebraic
formula for $\hh$ in terms of the structural matrix $\strmat$?
\end{problem}

\begin{problem}[Ihara--Selberg zeta function of EEAs]
\label{op:zeta}
Tian \cite[Section~6.2.3]{Tian2008} notes a connection between
evolution algebras and the Ihara--Selberg zeta function.  For an EEA
$\EA$, the Ihara zeta function $Z(\EA,u)=\prod_{[C]}(1-u^{|C|})^{-1}$
(product over prime cycles in $\varGamma(\EA,\natbasis)$) is known to
satisfy a Riemann Hypothesis if $\EA$ is a Ramanujan evolution algebra.
Develop the analogy between the Riemann Hypothesis for EEAs and the
classical Riemann Hypothesis.
\end{problem}

\begin{problem}[EEAs and 3-manifolds]
\label{op:3manifold}
Following Tian's programme \cite[Section~6.2.6]{Tian2008}, a
triangulation of a 3-manifold $M$ defines an evolution algebra
$\mathcal{A}(M,t)$.  When is the family $(\mathcal{A}(M,t_k))$ of
evolution algebras (over successive barycentric subdivisions $t_k$) an
EEA family?  Is this related to geometric or topological properties of
$M$ (e.g., hyperbolicity)?
\end{problem}

\begin{problem}[Automorphisms and symmetry of EEAs]
\label{op:aut}
For a Ramanujan evolution algebra $\EA$, compute the automorphism group
$\Aut(\EA)$ (as described by Elduque--Labra \cite{ElduqueLabra2019}).
Are Ramanujan EEAs more ``rigid'' (smaller automorphism group) than
non-Ramanujan EEAs of the same dimension and degree?
\end{problem}

\begin{problem}[Continuous EEAs]
\label{op:continuous}
Tian \cite[Section~6.2.4]{Tian2008} proposes continuous evolution
algebras $e_i\cdot e_i=\sum_j a_{ij}(t)e_j$ with time-dependent
structural constants.  Define a \emph{continuous EEA} in which
$\hh(\EA,\natbasis,t)\ge h$ for all $t$, and develop an analogue of
the mixing time result (Theorem~\ref{thm:mixing}) for continuous-time
Markov chains.
\end{problem}

\begin{problem}[Higher-dimensional analogues]
\label{op:higher}
Lubotzky \cite{Lubotzky2014} surveys higher-dimensional expander
complexes (coboundary expanders, cosystolic expanders, spectral
expanders).  Define and study \emph{higher-order EEAs} based on
simplicial complexes, and investigate which results of the present
paper generalise to this setting.
\end{problem}

% ==================================================================
% Bibliography
% ==================================================================

\begin{thebibliography}{99}

\bibitem{AlonMilman1985}
N.\ Alon and V.D.\ Milman,
$\lambda_1$, isoperimetric inequalities for graphs, and superconcentrators,
\emph{J.\ Combin.\ Theory Ser.\ B} \textbf{38} (1985), no.\ 1, 73--88.

\bibitem{BourgainGamburd2008}
J.\ Bourgain and A.\ Gamburd,
Uniform expansion bounds for Cayley graphs of $\mathrm{SL}_2(\mathbb{F}_p)$,
\emph{Ann.\ of Math.\ (2)} \textbf{167} (2008), no.\ 2, 625--642.

\bibitem{CadavidRodriguezRodini2019}
P.\ Cadavid, M.L.\ Rodi\~{n}o~Montoya, and P.M.\ Rodr\'{i}guez,
On the connection between evolution algebras, random walks and graphs,
\emph{J.\ Algebra Appl.} \textbf{19} (2020), no.\ 2, 2050023 (23 pp.).
\texttt{arXiv:1707.05897}.

\bibitem{CamachoGomezOmirovTurdibaev2010}
L.M.\ Camacho, J.R.\ G\'{o}mez, B.A.\ Omirov, and R.M.\ Turdibaev,
Some properties of evolution algebras,
\emph{Bull.\ Korean Math.\ Soc.} \textbf{50} (2013), no.\ 5, 1567--1586.
\texttt{arXiv:1004.1987}.

\bibitem{CasasLadraOmirovRozikov2010}
J.M.\ Casas, M.\ Ladra, B.A.\ Omirov, and U.A.\ Rozikov,
On evolution algebras,
\emph{Algebra Colloq.} \textbf{21} (2014), no.\ 2, 331--342.
\texttt{arXiv:1004.1050}.

\bibitem{ElduqueLabra2013}
A.\ Elduque and A.\ Labra,
Evolution algebras and graphs,
\emph{J.\ Algebra Appl.} \textbf{14} (2015), no.\ 7, 1550103 (10 pp.).
\texttt{arXiv:1312.4685}.

\bibitem{ElduqueLabra2019}
A.\ Elduque and A.\ Labra,
Evolution algebras, automorphisms, and graphs,
\emph{Linear Multilinear Algebra} \textbf{69} (2021), no.\ 2, 1--14.
\texttt{arXiv:1902.02191}.

\bibitem{Kowalski2019}
E.\ Kowalski,
\emph{An Introduction to Expander Graphs},
Cours Sp\'{e}cialis\'{e}s, Vol.\ 26,
Soci\'{e}t\'{e} Math\'{e}matique de France, Paris, 2019.

\bibitem{LPS1988}
A.\ Lubotzky, R.\ Phillips, and P.\ Sarnak,
Ramanujan graphs,
\emph{Combinatorica} \textbf{8} (1988), no.\ 3, 261--277.

\bibitem{Lubotzky2014}
A.\ Lubotzky,
Ramanujan complexes and high dimensional expanders,
\emph{Jpn.\ J.\ Math.} \textbf{9} (2014), no.\ 2, 137--169.

\bibitem{MSS2015}
A.W.\ Marcus, D.A.\ Spielman, and N.\ Srivastava,
Interlacing families II: Mixed characteristic polynomials and the Kadison--Singer problem,
\emph{Ann.\ of Math.\ (2)} \textbf{182} (2015), no.\ 1, 327--350.

\bibitem{Nilli1991}
A.\ Nilli,
On the second eigenvalue of a graph,
\emph{Discrete Math.} \textbf{91} (1991), no.\ 2, 207--210.

\bibitem{Pinsker1973}
M.S.\ Pinsker,
On the complexity of a concentrator,
in: \emph{7th International Teletraffic Conference}, Stockholm, 1973, pp.\ 318/1--318/4.

\bibitem{Tian2008}
J.P.\ Tian,
\emph{Evolution Algebras and Their Applications},
Lecture Notes in Mathematics, Vol.\ 1921,
Springer-Verlag, Berlin, 2008.

\bibitem{TianVojtechovsky2006}
J.P.\ Tian and P.\ Vojtechovsky,
Mathematical concepts of evolution algebras in non-Mendelian genetics,
\emph{Quasigroups Related Systems} \textbf{14} (2006), no.\ 1, 111--122.

\end{thebibliography}

\end{document}
